% ====================================================================
%  Global Cargo Ship Network Analysis
%  Authors: Aditya Gupta (2022031) & Debjit Banerji (2022146)
%  IIIT Delhi — Network Science Course Project, April 2026
%  Style: two-column, modelled on J. R. Soc. Interface (Kaluza 2010)
% ====================================================================
\documentclass[a4paper,10pt,twocolumn]{article}

% ── Fonts & encoding ─────────────────────────────────────────────────
\usepackage[T1]{fontenc}
\usepackage[utf8]{inputenc}
\usepackage{mathptmx}           % Times Roman body + math
\usepackage[scaled=0.90]{helvet}
\usepackage{microtype}

% ── Page geometry ────────────────────────────────────────────────────
\usepackage[a4paper,
            top=2.2cm, bottom=2.5cm,
            left=1.8cm, right=1.8cm,
            columnsep=0.65cm]{geometry}

% ── Colors ───────────────────────────────────────────────────────────
\usepackage[dvipsnames,table]{xcolor}
\definecolor{NavyHead}{RGB}{15,52,96}
\definecolor{TealLine}{RGB}{0,119,182}
\definecolor{TableH}{RGB}{218,232,247}

% ── Math ─────────────────────────────────────────────────────────────
\usepackage{amsmath,amssymb}

% ── Graphics ─────────────────────────────────────────────────────────
\usepackage{graphicx}
\usepackage{subcaption}
% stfloats allows figure*[b]; must load before float
\usepackage{stfloats}
\usepackage{float}
\captionsetup{font=small,labelfont=bf,labelsep=period,
              skip=5pt,width=\linewidth}
\captionsetup[subfigure]{font=small,labelfont=normalfont,
              labelformat=simple,labelsep=space}
\renewcommand\thesubfigure{(\alph{subfigure})}

% ── Tables ───────────────────────────────────────────────────────────
\usepackage{booktabs}
\usepackage{multirow}
\usepackage{array}
\usepackage{tabularx}
\usepackage{colortbl}
\newcolumntype{C}{>{\centering\arraybackslash}X}
\newcolumntype{R}{>{\raggedleft\arraybackslash}X}

% ── Section headings (Interface style) ───────────────────────────────
\usepackage{titlesec}
\titleformat{\section}
  {\normalfont\fontsize{9}{11}\bfseries\selectfont}
  {\thesection.}{0.4em}{\MakeUppercase}
\titleformat{\subsection}
  {\normalfont\fontsize{9}{11}\itshape\selectfont}
  {\thesubsection.}{0.4em}{}
\titleformat{\subsubsection}
  {\normalfont\fontsize{9}{11}\itshape\selectfont}
  {\thesubsubsection.}{0.4em}{}
\titlespacing*{\section}{0pt}{9pt plus 2pt minus 1pt}{3pt}
\titlespacing*{\subsection}{0pt}{7pt plus 1pt}{2pt}
\titlespacing*{\subsubsection}{0pt}{5pt plus 1pt}{2pt}

% ── Header / footer ──────────────────────────────────────────────────
\usepackage{fancyhdr}
\pagestyle{fancy}
\fancyhf{}
\fancyhead[L]{\fontsize{8}{9}\selectfont\itshape
  Complex network of global cargo ship movements}
\fancyhead[R]{\fontsize{8}{9}\selectfont\itshape
  Gupta \& Banerji}
\fancyfoot[C]{\fontsize{8}{9}\selectfont\thepage}
\renewcommand{\headrulewidth}{0.4pt}
\setlength{\headheight}{12pt}

% ── References ───────────────────────────────────────────────────────
\usepackage[authoryear,round,sort&compress]{natbib}
\bibpunct{(}{)}{;}{a}{,}{,}
\setlength{\bibsep}{2pt}
\renewcommand{\bibfont}{\small}

% ── Hyperlinks ───────────────────────────────────────────────────────
\usepackage[colorlinks=true,
            linkcolor=TealLine,
            citecolor=NavyHead,
            urlcolor=TealLine,
            bookmarks=true]{hyperref}

% ── Float parameters for two-column ─────────────────────────────────
\setlength{\dbltextfloatsep}{8pt plus 2pt minus 2pt}
\setlength{\dblfloatsep}{6pt plus 2pt minus 2pt}
\renewcommand{\dbltopfraction}{0.92}
\renewcommand{\textfraction}{0.05}
\renewcommand{\floatpagefraction}{0.80}
\renewcommand{\dblfloatpagefraction}{0.80}

% ── Misc ─────────────────────────────────────────────────────────────
\usepackage{enumitem}
\setlist[itemize]{leftmargin=1.2em,itemsep=1pt,parsep=0pt,topsep=2pt}
\setlist[enumerate]{leftmargin=1.4em,itemsep=1pt,parsep=0pt,topsep=2pt}
\usepackage{parskip}
\setlength{\parskip}{3pt}
\setlength{\parindent}{0pt}

% ── Full-width strip for title/abstract ──────────────────────────────
\usepackage{cuted}   % \strip environment

% ====================================================================
\begin{document}
\thispagestyle{empty}
% ====================================================================

% ── Full-width header banner ─────────────────────────────────────────
\begin{strip}
  \vspace*{-6pt}
  {\fontsize{7.5}{9}\selectfont
    \textit{Network Science Course Project, IIIT Delhi} \hfill
    \textit{April 2026}}\\[-3pt]
  \rule{\linewidth}{0.5pt}
  \vspace{2pt}
\end{strip}

% ── Full-width title block ────────────────────────────────────────────
\twocolumn[{%
\begin{@twocolumnfalse}
  \begin{center}
    {\fontsize{17}{21}\bfseries\selectfont
      The complex network of global cargo ship\\[3pt]
      movements: a 2015--2025 replication and extension\par}
    \vspace{9pt}
    {\fontsize{10.5}{13}\bfseries\selectfont
      Aditya Gupta and Debjit Banerji\par}
    \vspace{3pt}
    {\fontsize{9}{11}\itshape\selectfont
      Department of Computer Science and Engineering,
      Indraprastha Institute of Information Technology Delhi
      (IIIT-D),\\
      Okhla Phase III, New Delhi 110020, India\\
      Roll Nos: 2022031 (A. Gupta) \& 2022146 (D. Banerji)\par}
    \vspace{10pt}
  \end{center}
  %
  % Abstract
  {\fontsize{9}{12}\selectfont
  We replicate and substantially extend the landmark study of
  \citet{Kaluza2010} on the global cargo ship network (GCSN)
  using open vessel-tracking data from the Global Fishing Watch
  (GFW) API for two temporal snapshots---2015 and 2025.
  Our port-to-port networks comprise 7\,097 and 7\,717 ports
  connected by up to 188\,545 directed edges encoding individual
  voyage counts.
  We characterise both networks as \emph{small-world}
  (clustering $C\approx0.52$, average path length $L\approx3.3$,
  small-world coefficient $\sigma\approx66$--73) with exponentially
  distributed degree sequences and superlinear strength--degree
  scaling ($\alpha\approx1.4$), closely replicating
  Kaluza~et~al.'s findings despite an entirely different data source.
  Comparing 2015 to 2025, total voyages \emph{increased} by 9.7\%
  while the number of routes fell 3.3\% and mean degree fell 12.0\%,
  revealing a traffic-concentration paradox consistent with
  mega-vessel deployment.
  Modularity increased ($Q:0.699\to0.720$) and Louvain detection
  reveals 28--34 geographically coherent communities,
  50\% of which are stable over the decade.
  Triadic census analysis confirms strong over-representation of
  clique-like motifs.
  We add PageRank and k-core centrality, port trade specialisation,
  strategic chokepoint disruption, and a country-level network
  aggregation.
  Shanghai is the most critical single chokepoint;
  China dominates the country-level network in both years.\par}
  \vspace{6pt}
  {\fontsize{9}{11}\selectfont
  \textbf{Keywords:}\ complex network; cargo ships;
  small-world network; community detection;
  temporal analysis; network resilience\par}
  \vspace{12pt}
  \rule{\linewidth}{0.4pt}
  \vspace{6pt}
\end{@twocolumnfalse}
}]

% ====================================================================
\section{Introduction}
% ====================================================================

Maritime shipping carries over 80\,per cent of world trade by volume
\citep{UNCTAD2023}.
The global cargo ship network (GCSN)---where nodes are ports and
directed weighted edges are voyage counts---is one of Earth's most
critical infrastructure systems, central to logistics, pandemic
preparedness, and supply-chain risk management.

By 2006, 7.4 billion tonnes of goods were handled at the world's
ports \citep{IMO2006}, and the GCSN also acts as a vector for
invasive species through ballast water and hull fouling
\citep{Kaluza2010}.

\citet{Kaluza2010} analysed Lloyd's Register data from 2007 and
characterised the GCSN over 951 ports as a small-world network
\citep{Watts1998} with rich community structure and superlinear
strength--degree scaling $\langle s(k)\rangle\!\sim\!k^{1.36}$.
\citet{Hu2009} independently analysed the worldwide maritime
transportation network using Lloyd's data and reached similar
structural conclusions, further motivating the present replication.
Since then, shipping volumes have grown substantially and
disruptions---the COVID-19 pandemic (2020), the Russia--Ukraine
conflict (2022), and the 2021 Suez Canal blockage---have reshaped
routes.

This project addresses two questions:
(i)~\textit{Replication}---do the structural properties of
\citet{Kaluza2010} hold for GFW-derived data in 2015?
(ii)~\textit{Evolution}---how has topology, resilience and community
structure changed from 2015 to 2025?
Beyond replication, five extensions are presented: PageRank and
k-core centrality, port trade specialisation, strategic chokepoint
disruption, comprehensive temporal analysis, and a country-level
aggregation inspired by \citet{Ducruet2022} and \citet{Liu2025}.

% ====================================================================
\section{Data}
% ====================================================================

Port-visit events were downloaded from the \textbf{Global Fishing
Watch (GFW) API} for vessel types \texttt{CARGO} and
\texttt{CONTAINER} for calendar years 2015 and 2025.
For each vessel, events were sorted by entry time, consecutive
duplicate ports removed (suppress trivial self-loops), and
sequential port pairs extracted to form a directed weighted edge list
(source, destination, voyage~count).

\begin{table}[t]
\centering
\caption{Dataset overview.}
\label{tab:data}
{\small
\begin{tabular}{@{}lrr@{}}
\toprule
\rowcolor{TableH}
\textbf{Year} & \textbf{Unique ports} & \textbf{Directed edges} \\
\midrule
2015 & 7,097 & 188,545 \\
2025 & 7,717 & 182,300 \\
\bottomrule
\end{tabular}}
\end{table}

The Automatic Identification System (AIS), adopted from 2012, enables
GFW to aggregate real-time vessel positions into port-visit events
globally---the open-data equivalent of Lloyd's Register.

% ====================================================================
\section{Network construction}
% ====================================================================

From each edge list, four graph variants are derived:
$G_d$ (directed unweighted), $G_u$ (undirected unweighted),
$G_d^w$ (directed, edge weight = voyage count), and
$G_u^w$ (undirected, weight = sum of both directions).
All analysis uses NetworkX \citep{Hagberg2008}.
The directed networks have reciprocity $\approx\!0.45$: $\sim$59\%
of linked port pairs have edges in both directions, confirming
most routes operate bi-directionally though not with equal frequency.

% ====================================================================
\section{Unweighted structural analysis}
% ====================================================================

\subsection{Summary statistics}

Table~\ref{tab:summary} reports key metrics alongside the
\citet{Kaluza2010} baseline.
Both years exhibit a giant connected component (GCC) covering
$>99.5\%$ of all ports.
The clustering coefficient $C=0.522$ (2015) closely replicates
the paper's $C=0.49$.
The longer path length ($L=3.26$ vs.\ $2.50$) reflects the much
larger port universe (7\,097 vs.\ 951 ports): seven times more ports
increase absolute hop distances even when routing efficiency is
similar.
The small-world coefficient $\sigma\approx66$--73 confirms clustering
exceeds a random null by two orders of magnitude while path length
is only marginally longer.

\begin{table*}[t]
\centering
\caption{Unweighted network summary: 2015, 2025, and the
  \citeauthor{Kaluza2010} 2007 baseline.}
\label{tab:summary}
{\small
\begin{tabularx}{\textwidth}{@{}lCCC@{}}
\toprule
\rowcolor{TableH}
\textbf{Metric} & \textbf{2015} & \textbf{2025} & \textbf{2007\,$^\dagger$} \\
\midrule
Nodes $N$                              & 7\,097   & 7\,717   & 951    \\
Directed edges                         & 188\,545 & 182\,300 & ---    \\
Undirected edges                       & 145\,745 & 139\,499 & ---    \\
GCC fraction                           & 99.8\%   & 99.6\%   & ---    \\
Mean degree $\langle k\rangle$         & 41.07    & 36.15    & 44.61  \\
Max degree                             & 1\,233   & 1\,171   & ---    \\
Mean in-degree                         & 26.57    & 23.62    & ---    \\
Clustering $C$                         & 0.5223   & 0.4884   & 0.490  \\
$C(k)$ decay exponent $\beta$          & 0.296    & 0.273    & ---    \\
Avg path length $L$                    & 3.26     & 3.57     & 2.50   \\
Diameter (sampled)                     & 10       & 14       & ---    \\
Assortativity $r$                      & $-$0.069 & $-$0.008 & ---    \\
KNN slope $\nu$                        & $-$0.014 & $+$0.055 & ---    \\
Degree exponent $\mu$ (MLE)            & 1.306    & 1.314    & 1.710  \\
Small-world $\sigma$                   & 66.1     & 73.0     & ---    \\
Small-world $\gamma$\;($C/C_\text{rand}$) & 90.2  & 104.2    & ---    \\
\bottomrule
\multicolumn{4}{@{}p{\textwidth}@{}}{\footnotesize
  $^\dagger$\citet{Kaluza2010}, combined cargo network.}
\end{tabularx}}
\end{table*}

\subsection{Degree distribution}

\begin{figure*}[t]
  \centering
  \includegraphics[width=\linewidth]{Figures/01_degree_distribution.png}
  \caption{Degree distribution $P(k)$ for 2015 and 2025 in log--log
    (top row) and log-linear (bottom row) space, with MLE power-law
    (red dashed) and OLS exponential (blue solid) fits.}
  \label{fig:degree}
\end{figure*}

Figure~\ref{fig:degree} shows that $P(k)$ in log-linear space is
well-described by a straight line, confirming an exponential rather
than a power-law distribution.
The MLE exponent $\mu\approx1.31$ yields $R^2<0.85$ in log-log space
whereas the exponential achieves $R^2>0.95$ in log-linear space.
MLE fitting follows \citet{Clauset2009}.
This is fully consistent with \citet{Kaluza2010}'s conclusion that
the GCSN is \textit{not} scale-free: unlike social networks or the
internet, shipping routes are constrained by physical geography and
commercial economics that impose a characteristic connectivity scale.

\subsection{Clustering and hierarchical structure}

\begin{figure}[t]
  \centering
  \includegraphics[width=\linewidth]{Figures/01_clustering.png}
  \caption{Clustering coefficient $C(k)$ vs.\ degree $k$ (log--log)
    for 2015 and 2025.
    Negative slope $\beta\approx0.29$--$0.30$ confirms hierarchical
    organisation \citep{Ravasz2002}.}
  \label{fig:clustering}
\end{figure}

The decaying $C(k)$ profile (figure~\ref{fig:clustering}) with
$\beta\approx0.29$--$0.30$ confirms hierarchical organisation in both
years: peripheral ports form dense local clusters while high-degree
hubs bridge communities with sparser local neighbourhoods
\citep{Ravasz2002}.
The slight drop from $C=0.522$ to $0.488$ is consistent with
peripheral route loss---when low-degree ports lose links, their
triangular connections disappear proportionally more.

\subsection{Path length distribution}

\begin{figure}[t]
  \centering
  \includegraphics[width=\linewidth]{Figures/01_path_length.png}
  \caption{Shortest-path length distributions sampled from the GCC
    (200 random source nodes).
    Both distributions peak at $L=3$; the 2025 distribution is
    shifted right (mean 3.57 vs.\ 3.26), reflecting sparser
    connectivity.}
  \label{fig:pathlen}
\end{figure}

Both distributions peak at $L=3$ (figure~\ref{fig:pathlen}),
confirming the small-world property: almost any two ports in the
world are connected within three intermediate stops.
The rightward shift of the 2025 distribution is a direct consequence
of the $-12\%$ decline in mean degree: fewer edges means slightly
longer routes, even though the network retains its overall
small-world character.

\subsection{Degree--degree correlations}

\begin{figure}[t]
  \centering
  \includegraphics[width=\linewidth]{Figures/01_knn.png}
  \caption{Average nearest-neighbour degree $k_\text{nn}(k)$
    vs.\ $k$ (log--log).
    The slope $\nu$ changes sign from $-0.014$ (2015) to $+0.055$
    (2025), indicating a shift toward assortative mixing.}
  \label{fig:knn}
\end{figure}

The KNN profile (figure~\ref{fig:knn}) reveals a sign change in the
degree-mixing slope between 2015 ($\nu=-0.014$, weak
disassortativity) and 2025 ($\nu=+0.055$, mild assortativity).
This shift---and the accompanying assortativity change from
$r=-0.069$ to $r=-0.008$ \citep{Newman2002}---suggests that by 2025
the network periphery has thinned, leaving mainly hub-to-hub
connections that drive the network toward assortative mixing.

\subsection{Rich-club coefficient}

\begin{figure}[t]
  \centering
  \includegraphics[width=\linewidth]{Figures/01_rich_club.png}
  \caption{Rich-club coefficient $\phi(k)$ vs.\ degree $k$.
    The sharp rise for $k>200$ indicates a strong rich-club
    effect: the highest-degree ports form a dense mutual core.}
  \label{fig:richclub}
\end{figure}

The sharp rise in $\phi(k)$ above $k=200$ (figure~\ref{fig:richclub})
confirms a strong rich-club effect: top-tier hubs such as Shanghai,
Rotterdam, and Singapore are densely interconnected, forming the
backbone of the global shipping system.
This hub-and-spoke architecture is the structural complement of the
hierarchical clustering pattern observed in $C(k)$.

\subsection{Top ports by degree}

\begin{figure*}[t]
  \centering
  \begin{subfigure}[b]{0.48\linewidth}
    \includegraphics[width=\linewidth]{Figures/01_top_ports_2015.png}
    \caption{2015}
  \end{subfigure}
  \hfill
  \begin{subfigure}[b]{0.48\linewidth}
    \includegraphics[width=\linewidth]{Figures/01_top_ports_2025.png}
    \caption{2025}
  \end{subfigure}
  \caption{Top-25 ports by undirected degree for (\textit{a})~2015
    and (\textit{b})~2025.}
  \label{fig:topports}
\end{figure*}

Shanghai, Rotterdam--Maasvlakte, and Antwerp consistently occupy the
top tier in both years.
A notable shift is Pengerang (Malaysia) dropping sharply from a
high 2015 rank by 2025---linked to downstream refinery schedule
changes in Southeast Asia.
Conversely, St.\,Petersburg and Skagen rise, reflecting increased
Baltic Sea traffic.
The 2025 ranking is moderately more compressed (max degree
$1171$ vs.\ $1233$), consistent with overall route consolidation.

% ====================================================================
\section{Weighted analysis}
% ====================================================================

\subsection{Link weight and strength distributions}

\begin{figure*}[t]
  \centering
  \begin{subfigure}[b]{0.48\linewidth}
    \includegraphics[width=\linewidth]{Figures/02_weight_distribution.png}
    \caption{Link weight $P(w)$}
  \end{subfigure}
  \hfill
  \begin{subfigure}[b]{0.48\linewidth}
    \includegraphics[width=\linewidth]{Figures/02_strength_distribution.png}
    \caption{Node strength $P(s)$}
  \end{subfigure}
  \caption{Heavy-tailed distributions of (\textit{a})~link weight
    and (\textit{b})~node strength in log--log space.
    MLE fits give $\mu\approx1.64$ and $\eta\approx1.24$.}
  \label{fig:weightdist}
\end{figure*}

Both distributions are heavy-tailed (figure~\ref{fig:weightdist}):
a few extremely busy routes and a handful of dominant hub ports
account for a disproportionate share of total traffic.
The weight exponent $\mu=1.643$ (2015) lies within the range
$1.42$--$1.93$ reported by \citet{Kaluza2010} for individual cargo
sub-types.
The strength exponent $\eta=1.236$ is slightly higher than the
paper's $1.02$--$1.05$, attributable to the larger port universe
in our dataset.

\subsection{Kaluza et al.\ figure 2 replication}

\begin{figure*}[t]
  \centering
  \includegraphics[width=\linewidth]{Figures/02_fig2_composite_replication.png}
  \caption{Replication of \citet{Kaluza2010} figure~2b--d.
    (\textit{a})~Link-weight distribution $P(w)$.
    (\textit{b})~Node-strength distribution $P(s)$.
    (\textit{c})~Mean strength per degree $\langle s(k)\rangle\sim k^\alpha$.
    Steelblue = 2015; orange = 2025.}
  \label{fig:composite}
\end{figure*}

Figure~\ref{fig:composite} and table~\ref{tab:weighted} constitute a
full replication of Kaluza et~al.'s Table~1.
The strength--degree scaling exponent $\alpha=1.435$ (2015) and
$\alpha=1.346$ (2025) bracket the paper's $\alpha=1.36$:
highly-connected ports carry disproportionately more traffic.
The excellent fit ($R^2>0.98$) confirms superlinear traffic
concentration at network hubs in both years.

\begin{table*}[t]
\centering
\caption{Weighted network metrics: replication of
  \citeauthor{Kaluza2010} Table~1.}
\label{tab:weighted}
{\small
\begin{tabularx}{\textwidth}{@{}lCCC@{}}
\toprule
\rowcolor{TableH}
\textbf{Metric} & \textbf{2015} & \textbf{2025} & \textbf{2007\,$^\dagger$} \\
\midrule
Total voyages                              & 1\,791\,795 & 1\,965\,325 & ---  \\
Mean journeys/link $\langle J\rangle$      & 9.50        & 10.78       & 10.4 \\
Max link weight                            & 10\,269     & 12\,779     & ---  \\
Mean strength $\langle s\rangle$           & 504.9       & 509.3       & ---  \\
Weight exponent $\mu$                      & 1.643       & 1.617       & $\approx$1.42 \\
Strength exponent $\eta$                   & 1.236       & 1.233       & 1.02 \\
Strength--degree exponent $\alpha$         & 1.435       & 1.346       & 1.36 \\
$s$--$k$ scaling $R^2$                     & 0.993       & 0.983       & ---  \\
Weighted clustering $C_w$                  & 0.0002      & 0.0001      & ---  \\
\bottomrule
\multicolumn{4}{@{}p{\textwidth}@{}}{\footnotesize
  $^\dagger$\citet{Kaluza2010}, combined cargo network.}
\end{tabularx}}
\end{table*}

\subsection{Weighted clustering and KNN}

\begin{figure*}[t]
  \centering
  \begin{subfigure}[b]{0.48\linewidth}
    \includegraphics[width=\linewidth]{Figures/02_weighted_clustering.png}
    \caption{Weighted clustering $C_w(k)$}
  \end{subfigure}
  \hfill
  \begin{subfigure}[b]{0.48\linewidth}
    \includegraphics[width=\linewidth]{Figures/02_weighted_knn.png}
    \caption{Weighted KNN $k_\text{nn}^w(k)$}
  \end{subfigure}
  \caption{Weighted structural profiles \citep{Barrat2004}.
    (\textit{a})~Weighted clustering is two orders of magnitude below
    unweighted $C(k)$: triangles are structurally present but
    most traffic concentrates on a few heavy routes.
    (\textit{b})~Weighted KNN declines with $k$: hubs' heavy
    connections go to lower-traffic ports.}
  \label{fig:wclustering}
\end{figure*}

The large gap between unweighted ($C\approx0.52$) and weighted
($C_w\approx0.0002$) clustering (figure~\ref{fig:wclustering}a)
is a hallmark of hub-and-spoke structure: many triangles exist
topologically but the actual traffic through them is negligible.
Most voyage volume concentrates on the handful of routes connecting
major container hubs, leaving the majority of triangles
effectively ``dark'' in traffic terms.

\subsection{Top ports by traffic strength}

\begin{figure*}[t]
  \centering
  \begin{subfigure}[b]{0.48\linewidth}
    \includegraphics[width=\linewidth]{Figures/02_top_ports_strength_2015.png}
    \caption{2015}
  \end{subfigure}
  \hfill
  \begin{subfigure}[b]{0.48\linewidth}
    \includegraphics[width=\linewidth]{Figures/02_top_ports_strength_2025.png}
    \caption{2025}
  \end{subfigure}
  \caption{Top-25 ports by node strength (total voyage count) for
    (\textit{a})~2015 and (\textit{b})~2025.}
  \label{fig:topstrength}
\end{figure*}

Shanghai and Rotterdam--Maasvlakte dominate in both years
(figure~\ref{fig:topstrength}).
Rank--strength correlation $\rho\approx0.91$ between degree and
strength confirms that high-degree hubs also carry the most
aggregate traffic.
Total voyages grew $+9.7\%$ to 1.97 million in 2025, but this
growth is unevenly distributed: the top-10 hubs by strength captured
a disproportionate share of new traffic, deepening inequality.

\subsection{Most active bilateral routes}

\begin{figure*}[t]
  \centering
  \begin{subfigure}[b]{0.48\linewidth}
    \includegraphics[width=\linewidth]{Figures/02_top_routes_2015.png}
    \caption{2015}
  \end{subfigure}
  \hfill
  \begin{subfigure}[b]{0.48\linewidth}
    \includegraphics[width=\linewidth]{Figures/02_top_routes_2025.png}
    \caption{2025}
  \end{subfigure}
  \caption{Top-30 bilateral routes by voyage count for
    (\textit{a})~2015 and (\textit{b})~2025.
    Max link weight grew from 10\,269 to 12\,779 ($+24.4\%$).}
  \label{fig:toproutes}
\end{figure*}

The heaviest routes are overwhelmingly intra-regional---North Sea,
East Asia, and the Mediterranean---with a handful of inter-regional
links connecting Chinese and European hub clusters
(figure~\ref{fig:toproutes}).
The maximum link weight grew $+24.4\%$ to 12\,779 voyages by 2025,
reflecting dedicated mega-vessel services on a few arterial
corridors.
The composition of the top-30 routes is highly stable: route rank
correlation $\rho\approx0.87$ between years.

\subsection{Port trade specialisation}

\begin{figure*}[t]
  \centering
  \begin{subfigure}[b]{0.48\linewidth}
    \includegraphics[width=\linewidth]{Figures/02_port_specialisation_2015.png}
    \caption{2015}
  \end{subfigure}
  \hfill
  \begin{subfigure}[b]{0.48\linewidth}
    \includegraphics[width=\linewidth]{Figures/02_port_specialisation_2025.png}
    \caption{2025}
  \end{subfigure}
  \caption{Port trade specialisation: in-strength vs.\ out-strength
    (log--log) coloured by asymmetry index
    $A\!=\!(s_\text{out}\!-\!s_\text{in})/(s_\text{out}\!+\!s_\text{in})$.
    Blue = net importer; red = net exporter.
    Bar charts show the most asymmetric ports by $|A|$.}
  \label{fig:specialisation}
\end{figure*}

The vast majority of high-strength ports cluster tightly along the
diagonal ($|A|<0.002$), confirming they function as near-perfectly
balanced transshipment relays (figure~\ref{fig:specialisation}).
Asymmetric ports are predominantly low-volume, single-commodity
terminals---bulk mineral exporters or energy import facilities---
where directional traffic imbalance reflects their production or
consumption role rather than transit.
This pattern is stable between 2015 and 2025, suggesting port
trade roles are structurally locked-in at decadal timescales.

% ====================================================================
\section{Centrality analysis}
% ====================================================================

\subsection{Betweenness centrality}

\begin{figure*}[t]
  \centering
  \begin{subfigure}[b]{0.48\linewidth}
    \includegraphics[width=\linewidth]{Figures/01_betweenness_centrality_2015.png}
    \caption{2015: Shanghai leads ($\mathrm{BC}=0.075$)}
  \end{subfigure}
  \hfill
  \begin{subfigure}[b]{0.48\linewidth}
    \includegraphics[width=\linewidth]{Figures/01_betweenness_centrality_2025.png}
    \caption{2025: Rotterdam overtakes ($\mathrm{BC}=0.078$)}
  \end{subfigure}
  \caption{Top-25 ports by betweenness centrality (BC), approximated
    with $k=500$ random pivot nodes \citep{Brandes2001}.}
  \label{fig:bc}
\end{figure*}

Shanghai leads BC in 2015 ($0.075$), followed by
Rotterdam--Maasvlakte ($0.072$) and Antwerp ($0.056$)
(figure~\ref{fig:bc}).
By 2025, Rotterdam--Maasvlakte overtakes Shanghai ($0.078$ vs.\
$0.077$), reflecting sustained growth of European transshipment.
Ports such as Gibraltar and Skagen achieve disproportionately high
BC relative to their degree, a signature of geographic chokepoint
positions on major shipping lanes.
Spearman rank correlation between BC and degree is
$\rho\approx0.80$--$0.85$: BC captures routing bridgeness that
degree alone cannot.

\subsection{PageRank centrality}

\begin{figure*}[t]
  \centering
  \begin{subfigure}[b]{0.48\linewidth}
    \includegraphics[width=\linewidth]{Figures/01_pagerank_2015.png}
    \caption{2015: Shanghai ($4.48\times10^{-3}$)}
  \end{subfigure}
  \hfill
  \begin{subfigure}[b]{0.48\linewidth}
    \includegraphics[width=\linewidth]{Figures/01_pagerank_2025.png}
    \caption{2025: Shanghai ($3.84\times10^{-3}$)}
  \end{subfigure}
  \caption{Top-25 ports by PageRank (directed, unweighted,
    damping $\alpha=0.85$) \citep{Page1999}.}
  \label{fig:pagerank}
\end{figure*}

PageRank captures \textit{gateway importance}: a port fed by many
important hubs scores highly even without lying on many shortest
paths (figure~\ref{fig:pagerank}).
Shanghai remains the top PageRank port in both years, while Pengerang
(Malaysia) drops from 5th in 2015 to below the top-25 in 2025
($5.51\times10^{-4}$), reflecting a major Southeast Asian hub shift.
The three centrality measures are complementary: Spearman
$\rho(\text{PR, degree})\approx0.88$;
$\rho(\text{PR, BC})\approx0.72$;
$\rho(\text{BC, degree})\approx0.82$.

\subsection{k-core decomposition}

\begin{figure*}[t]
  \centering
  \includegraphics[width=\linewidth]{Figures/01_kcore_decomposition.png}
  \caption{k-core decomposition \citep{Seidman1983}.
    (\textit{a})~k-shell size distribution.
    (\textit{b})~Cumulative node fraction vs.\ $k$.
    (\textit{c})~Coreness vs.\ degree (2015).
    Maximum core $k_\text{max}=74$ (2015), $72$ (2025), dominated by
    East Asian ports.}
  \label{fig:kcore}
\end{figure*}

The innermost core ($k_\text{max}=74$) identifies the dense backbone:
ports include Tianjin, Taichung, Daesan (South Korea), Weihai, and
Shenzhen---all East Asian (figure~\ref{fig:kcore}).
This confirms the central role of intra-Asia container shipping in
the global network backbone.
The slight contraction to $k_\text{max}=72$ over the decade
signals a marginal weakening of the densest core, consistent with
route consolidation onto fewer mega-vessel services.
The power-law-like coreness--degree scatter in panel~(c) confirms
that high-degree nodes are necessary but not sufficient for high
coreness: geographic positioning in a dense local neighbourhood
also matters.

% ====================================================================
\section{Temporal analysis: 2015 vs.\ 2025}
% ====================================================================

\subsection{Structural change overview}

\begin{table*}[t]
\centering
\caption{Full structural comparison: 2015 vs.\ 2025.}
\label{tab:temporal}
{\small
\begin{tabularx}{\textwidth}{@{}lCCC@{}}
\toprule
\rowcolor{TableH}
\textbf{Metric} & \textbf{2015} & \textbf{2025} & \textbf{Change} \\
\midrule
\multicolumn{4}{@{}l@{}}{\footnotesize\itshape Topology} \\
Nodes $N$            & 7\,097      & 7\,717      & $+$8.7\%    \\
Directed edges       & 188\,545    & 182\,300    & $-$3.3\%    \\
Directed density     & 0.003744    & 0.003062    & $-$18.2\%   \\
Connected components & 8           & 5           & $-$37.5\%   \\
SCCs (directed)      & 219         & 203         & $-$7.3\%    \\[3pt]
\multicolumn{4}{@{}l@{}}{\footnotesize\itshape Degree \& connectivity} \\
Mean degree          & 41.07       & 36.15       & $-$12.0\%   \\
Clustering $C$       & 0.5223      & 0.4884      & $-$6.5\%    \\
Path length $L$      & 3.26        & 3.57        & $+$8.7\%    \\
Assortativity $r$    & $-$0.069    & $-$0.008    & $+$88.5\%   \\[3pt]
\multicolumn{4}{@{}l@{}}{\footnotesize\itshape Traffic} \\
Total voyages        & 1\,791\,795 & 1\,965\,325 & $+$9.7\%    \\
$\langle J\rangle$   & 9.50        & 10.78       & $+$13.4\%   \\
Max link weight      & 10\,269     & 12\,779     & $+$24.4\%   \\[3pt]
\multicolumn{4}{@{}l@{}}{\footnotesize\itshape Mesoscale} \\
Communities $c$      & 28          & 34          & $+$21.4\%   \\
Modularity $Q$       & 0.699       & 0.720       & $+$3.0\%    \\
\bottomrule
\end{tabularx}}
\end{table*}

Table~\ref{tab:temporal} reveals the \textit{traffic-concentration
paradox}: total voyages grew $+9.7\%$ while the number of routes
fell $-3.3\%$ and mean degree fell $-12.0\%$.
More ships are traversing fewer routes, each carrying heavier
traffic---directly consistent with the industry-wide trend of
mega-vessel deployment \citep{Notteboom2012,Liu2025}.

\subsection{Distribution shifts}

\begin{figure*}[t]
  \centering
  \includegraphics[width=\linewidth]{Figures/03_distribution_comparison.png}
  \caption{Complementary CDFs of degree, in-degree, and edge weight
    for 2015 (steelblue) and 2025 (orange).
    Kolmogorov--Smirnov tests: all shifts significant at
    $p<10^{-4}$.}
  \label{fig:distr}
\end{figure*}

All three distributions shift significantly between years
(figure~\ref{fig:distr}).
The degree CCDF shifts leftward (fewer connections per port) while
the weight CCDF shifts rightward (heavier individual routes).
These simultaneous shifts in opposite directions are the quantitative
signature of the traffic-concentration paradox: the network
reconfigures from many light routes to fewer heavy ones.

\subsection{Network overlap and edge Jaccard similarity}

\begin{figure}[t]
  \centering
  \includegraphics[width=\linewidth]{Figures/03_overlap.png}
  \caption{Network overlap between 2015 and 2025: Jaccard similarity
    of node and edge sets, and the fraction of traffic on
    persistent routes.}
  \label{fig:overlap}
\end{figure}

Despite the decade gap, a large core of routes persists between years
(figure~\ref{fig:overlap}).
The high node-set overlap ($>$85\%) confirms that the port universe
is essentially stable; it is the edge set that changes.
Critically, persistent routes carry the \textit{majority} of total
traffic, implying that the network's structural core---the
high-traffic arteries---is highly stable, while peripheral
low-traffic routes appear and disappear.

\subsection{Route weight changes}

\begin{figure}[t]
  \centering
  \includegraphics[width=\linewidth]{Figures/03_route_changes.png}
  \caption{Distribution of log-ratio
    $\log_2(w_{2025}/w_{2015})$ for persistent routes.
    Near-symmetric distribution indicates balanced intensification
    and reduction, but with increased variance.}
  \label{fig:routechanges}
\end{figure}

The log-ratio distribution for persistent routes is approximately
symmetric around zero (figure~\ref{fig:routechanges}): on average,
routes neither gain nor lose traffic.
However, the increased variance indicates polarisation---some routes
became dramatically busier while others nearly ceased.
This ``route bifurcation'' is consistent with shipper consolidation
onto fewer carriers on main corridors, while marginal routes lose
traffic to larger hub calls.

\subsection{Port importance evolution}

\begin{figure}[t]
  \centering
  \includegraphics[width=\linewidth]{Figures/03_port_evolution.png}
  \caption{Rank evolution of port strength for the top-50 ports.
    Hub rankings are remarkably stable ($\rho\approx0.93$),
    with notable risers (Ningbo, Piraeus) and fallers (Pengerang).}
  \label{fig:portevolution}
\end{figure}

The Spearman rank correlation $\rho\approx0.93$ between 2015 and
2025 port strength rankings (figure~\ref{fig:portevolution}) is
strikingly high over a decade of major geopolitical and economic
disruptions.
Notable risers include Ningbo-Zhoushan (now the world's busiest
port by tonnage) and Piraeus (Chinese investment through
COSCO Shipping).
Pengerang's steep decline reflects the scaling-back of the
Pengerang Integrated Petroleum Complex following the 2019 fire.

\subsection{New and lost ports}

\begin{figure}[t]
  \centering
  \includegraphics[width=\linewidth]{Figures/03_new_lost_ports.png}
  \caption{Characterisation of new ports (in 2025 only) and lost
    ports (in 2015 only).
    Both groups are dominated by low-degree, low-strength peripheral
    nodes consistent with marginal terminals entering/leaving
    service.}
  \label{fig:newlost}
\end{figure}

New and lost ports are overwhelmingly peripheral, low-degree nodes
(figure~\ref{fig:newlost}).
This confirms that the network's structural changes are driven by the
appearance and disappearance of marginal terminals---small coastal
facilities and single-berth terminals---rather than
any reorganisation of major hubs.
The 620 new ports in 2025 reflect expanding coverage in emerging
markets (Southeast Asia, West Africa); the lost ports reflect
terminal consolidation in mature markets.

\subsection{Traffic concentration}

\begin{figure*}[t]
  \centering
  \includegraphics[width=\linewidth]{Figures/03_lorenz_gini.png}
  \caption{Lorenz curves for edge weights (left), node strengths
    (centre), and degree (right).
    Gini coefficients exceed 0.85 for edge weights in both years.
    All three Gini values increased from 2015 to 2025.}
  \label{fig:lorenz}
\end{figure*}

Approximately 5\% of routes carry the majority of total voyage
traffic (figure~\ref{fig:lorenz}).
This extreme inequality---Gini $>0.85$ for edge weights---is
substantially more concentrated than income inequality in the most
unequal societies (Gini $\approx0.6$).
The 2015-to-2025 increase in all three Gini coefficients is the
clearest quantitative indicator of the traffic-concentration paradox.

\subsection{Reciprocity and route directionality}

\begin{figure}[t]
  \centering
  \includegraphics[width=\linewidth]{Figures/03_reciprocity.png}
  \caption{Reciprocity and directional asymmetry of bilateral routes.
    Reciprocity $\approx0.76$, stable across years.
    The asymmetry distribution is unimodal and centred near zero.}
  \label{fig:reciprocity}
\end{figure}

Network reciprocity is $\approx0.76$ and is nearly unchanged
($-0.003$) between years (figure~\ref{fig:reciprocity}).
This confirms the strongly bidirectional nature of cargo shipping:
almost all major routes operate in both directions with roughly
similar volumes.
The unimodal, near-symmetric asymmetry distribution shows that
strongly one-directional routes are rare, with the few exceptions
being dedicated tanker or bulk-cargo corridors.

\subsection{Network robustness}

\begin{figure}[t]
  \centering
  \includegraphics[width=\linewidth]{Figures/03_robustness.png}
  \caption{GCC fraction vs.\ removed-node fraction under targeted
    (highest-degree first) and random attack.
    Targeted attack causes rapid collapse; random removal has
    negligible effect.}
  \label{fig:robustness}
\end{figure}

The GCSN is highly robust to random failures but fragile under
targeted attack (figure~\ref{fig:robustness}): removing the top
15\% of degree hubs reduces the GCC below 50\%.
This robustness--vulnerability dichotomy is characteristic of
heterogeneous networks \citep{Barabasi1999} and has critical
implications for pandemic disease spread and supply-chain planning.
The 2025 network is marginally more vulnerable under targeted
attack, consistent with the lower mean degree.

\subsection{Chokepoint disruption analysis}

\begin{figure}[t]
  \centering
  \includegraphics[width=\linewidth]{Figures/03_chokepoint_disruption.png}
  \caption{Impact of removing individual strategic ports on
    (\textit{a})~GCC fraction reduction and
    (\textit{b})~APL increase.
    Shanghai causes the largest GCC reduction (0.113\%, 2015).}
  \label{fig:chokepoints}
\end{figure}

\begin{table}[t]
\centering
\caption{Chokepoint disruption: single-port removal impact.}
\label{tab:chokepoints}
{\small
\begin{tabular}{@{}lcccc@{}}
\toprule
\rowcolor{TableH}
 & \multicolumn{2}{c}{$\Delta$GCC (\%)} &
   \multicolumn{2}{c}{$\Delta\langle L\rangle$} \\
\cmidrule(lr){2-3}\cmidrule(lr){4-5}
Port & 2015 & 2025 & 2015 & 2025 \\
\midrule
Shanghai  & 0.113 & 0.065 & 0.044 & 0.022 \\
Rotterdam & 0.000 & 0.039 & 0.047 & 0.015 \\
Singapore & 0.014 & 0.026 & 0.030 & 0.007 \\
Busan     & 0.000 & 0.013 & 0.028 & 0.003 \\
Hamburg   & 0.028 & 0.013 & 0.017 & 0.004 \\
Port Said & 0.000 & 0.013 & 0.021 & 0.005 \\
\bottomrule
\end{tabular}}
\end{table}

No single port removal fragments the network
(table~\ref{tab:chokepoints}, figure~\ref{fig:chokepoints}):
the GCSN has no single point of failure for connectivity.
However, APL increases confirm significant efficiency degradation.
Rotterdam's APL impact ($\Delta L=0.047$) exceeds Shanghai's in 2015,
reflecting its critical gateway role for European inbound cargo.
Shanghai's declining GCC impact (halved to 0.065\% by 2025) suggests
the network adapted around it as Ningbo and other Chinese hubs grew.

\subsection{Hub stability}

\begin{figure}[t]
  \centering
  \includegraphics[width=\linewidth]{Figures/03_hub_stability.png}
  \caption{Spearman rank correlation of port strength (left) and
    degree (right) between years.
    $\rho\approx0.93$ for both measures: hub rankings are
    remarkably stable across the decade.}
  \label{fig:stability}
\end{figure}

The high rank stability ($\rho\approx0.93$, figure~\ref{fig:stability})
despite a decade of major disruptions is a striking result.
The COVID-19 pandemic (2020), US--China tariff wars, and the
Suez Canal blockage (2021) all generated short-term disruptions
in the news cycle, yet the structural hierarchy of ports is almost
entirely preserved.
This suggests that port infrastructure investments and shipping
alliances create a ``structural lock-in'' that resists rapid
reorganisation.

\subsection{Summary dashboard}

\begin{figure*}[t]
  \centering
  \includegraphics[width=\linewidth]{Figures/03_summary_dashboard.png}
  \caption{Visual summary dashboard: eight key metrics side-by-side
    for 2015 and 2025.
    The dashboard encapsulates the traffic-concentration paradox
    (more voyages, fewer routes), the robustness--vulnerability
    trade-off, and increasing modularity.}
  \label{fig:dashboard}
\end{figure*}

The summary dashboard (figure~\ref{fig:dashboard}) provides an
at-a-glance synthesis of the decade's structural evolution.
The simultaneous leftward shift in the degree distribution and
rightward shift in the weight distribution---visible in the same
panel---is the clearest visual representation of the
traffic-concentration paradox.

% ====================================================================
\section{Community detection and analysis}
% ====================================================================

\subsection{Detection method and summary}

Louvain community detection \citep{Blondel2008} was applied to
$G_u$ for both years.

\begin{table}[t]
\centering
\caption{Community detection summary.}
\label{tab:communities}
{\small
\begin{tabular}{@{}lccc@{}}
\toprule
\rowcolor{TableH}
\textbf{Metric} & \textbf{2015} & \textbf{2025} & \textbf{2007\,$^\dagger$} \\
\midrule
Algorithm           & Louv. & Louv. & L--N  \\
Communities $c$     & 28    & 34    & 12/7/6 \\
Modularity $Q$      & 0.699 & 0.720 & 0.605--0.716 \\
Largest (ports)     & 1\,990 & 1\,735 & --- \\
Intra-edge ratio    & 70.2\% & 67.6\% & --- \\
Intra-traffic ratio & 84.2\% & 83.4\% & --- \\
Coverage            & 0.702  & 0.676  & --- \\
Performance         & 0.860  & 0.894  & --- \\
Mean conductance    & 0.096  & 0.094  & --- \\
Stable ($J{\geq}0.25$) & 14/28 & ---  & --- \\
\bottomrule
\multicolumn{4}{@{}l@{}}{\footnotesize
  $^\dagger$\citet{Kaluza2010}; Louv.\ = Louvain;
  L--N = Leicht--Newman.}
\end{tabular}}
\end{table}

The community count grew from 28 to 34 ($+21\%$) and modularity
from $Q=0.699$ to $0.720$ (table~\ref{tab:communities}): regional
trade blocs tightened over the decade.
These values exceed the original paper's $Q=0.605$--$0.716$,
suggesting GFW combined data produces at least equally well-defined
communities.
Intra-community traffic remains near-constant at $\approx84\%$,
confirming that the vast majority of global shipping is regional.

\subsection{Community size and strength}

\begin{figure*}[t]
  \centering
  \begin{subfigure}[b]{0.48\linewidth}
    \includegraphics[width=\linewidth]{Figures/04_community_size.png}
    \caption{Community size (ports)}
  \end{subfigure}
  \hfill
  \begin{subfigure}[b]{0.48\linewidth}
    \includegraphics[width=\linewidth]{Figures/04_community_strength.png}
    \caption{Community traffic strength}
  \end{subfigure}
  \caption{(\textit{a})~Community size distribution: a single giant
    community coexists with many small peripheral ones.
    (\textit{b})~Community strength: the giant community captures
    a disproportionate share of all voyage traffic.}
  \label{fig:commsizes}
\end{figure*}

A single giant community ($\approx$1\,800--2\,000 ports) dominates
the size distribution (figure~\ref{fig:commsizes}a), surrounded by
a power-law-like tail of smaller communities.
The giant community's traffic share (panel~b) exceeds its port share,
confirming the ``mega-hub'' effect: the core community contains the
world's busiest ports and thus captures the majority of total traffic.

\subsection{Community structural profile}

\begin{figure}[t]
  \centering
  \includegraphics[width=\linewidth]{Figures/04_community_profile_scatter.png}
  \caption{Per-community profile: internal density vs.\ size,
    coloured by mean conductance.
    Smaller communities are denser and better-separated; the giant
    community is sparse with porous boundaries.}
  \label{fig:commprofile}
\end{figure}

The inverse relationship between community size and internal density
(figure~\ref{fig:commprofile}) is a standard feature of
hierarchically-structured networks.
Small communities (e.g.\ Baltic Sea cluster, Oceania cluster) are
highly cohesive geographic blocs.
The giant community is structurally porous---it contains ports from
multiple geographic regions connected through major transshipment
hubs, and correspondingly has lower internal density and higher
conductance.

\subsection{Inter-community traffic}

\begin{figure*}[t]
  \centering
  \begin{subfigure}[b]{0.48\linewidth}
    \includegraphics[width=\linewidth]{Figures/04_community_heatmap_2015.png}
    \caption{2015}
  \end{subfigure}
  \hfill
  \begin{subfigure}[b]{0.48\linewidth}
    \includegraphics[width=\linewidth]{Figures/04_community_heatmap_2025.png}
    \caption{2025}
  \end{subfigure}
  \caption{Inter-community traffic matrices (log scale) for
    (\textit{a})~2015 and (\textit{b})~2025.
    Diagonal dominance: 84\% of all traffic is intra-community.}
  \label{fig:heatmap}
\end{figure*}

The diagonal-dominant structure (figure~\ref{fig:heatmap}) confirms
strong intra-community flows: 84\% of all voyages remain within
the same community.
The dominant off-diagonal block connects the European
(Rotterdam-hub) and East-Asian (Shanghai-hub) communities,
representing the trans-Eurasian container trunk route.
The slight tightening of diagonal dominance in 2025 (despite more
communities) reflects the increasing modularity: while trade blocs
split into more sub-groups, each sub-group becomes more
self-contained.

\subsection{Community flow balance}

\begin{figure*}[t]
  \centering
  \begin{subfigure}[b]{0.48\linewidth}
    \includegraphics[width=\linewidth]{Figures/04_community_flow_2015.png}
    \caption{2015}
  \end{subfigure}
  \hfill
  \begin{subfigure}[b]{0.48\linewidth}
    \includegraphics[width=\linewidth]{Figures/04_community_flow_2025.png}
    \caption{2025}
  \end{subfigure}
  \caption{Community-level flow balance (net importer/exporter) for
    (\textit{a})~2015 and (\textit{b})~2025.
    Most communities are near-balanced; a few are specialised net
    exporters or importers.}
  \label{fig:commflow}
\end{figure*}

Most communities are near-balanced in directed voyage flow
(figure~\ref{fig:commflow}).
Specialised net-exporter communities tend to correspond to
resource-extraction regions (e.g.\ Australasian bulk commodities,
Persian Gulf energy corridors), while net-importer communities
reflect manufacturing consumption hubs.
The flow-balance landscape is broadly stable between 2015 and 2025,
consistent with the structural lock-in observed in port rankings.

\subsection{Community conductance and quality}

\begin{figure*}[t]
  \centering
  \begin{subfigure}[b]{0.48\linewidth}
    \includegraphics[width=\linewidth]{Figures/04_conductance_2015.png}
    \caption{2015: mean $\varphi=0.096$}
  \end{subfigure}
  \hfill
  \begin{subfigure}[b]{0.48\linewidth}
    \includegraphics[width=\linewidth]{Figures/04_conductance_2025.png}
    \caption{2025: mean $\varphi=0.094$}
  \end{subfigure}
  \caption{Community conductance distributions for (\textit{a})~2015
    and (\textit{b})~2025.
    Low conductance ($\varphi\approx0.09$) confirms well-separated
    communities with minimal cross-boundary traffic.}
  \label{fig:conductance}
\end{figure*}

The low mean conductance $\varphi\approx0.09$--$0.10$
(figure~\ref{fig:conductance}) confirms that Louvain communities are
well-separated: cross-boundary traffic is modest relative to
within-community traffic.
The slight decrease from $0.096$ to $0.094$ indicates marginally
better community isolation in 2025, consistent with the modularity
increase.
Communities with high conductance tend to be geographically
diffuse communities that bridge multiple trade regions.

\subsection{Community persistence}

\begin{figure}[t]
  \centering
  \includegraphics[width=\linewidth]{Figures/04_community_jaccard.png}
  \caption{Jaccard similarity matrix between 2015 and 2025 communities
    (top-20 sub-matrix).
    14/28 communities are stable ($J\geq0.25$).}
  \label{fig:jaccard}
\end{figure}

Fourteen of 28 communities (50\%) persist from 2015 to 2025 with
Jaccard similarity $J\geq0.25$ (figure~\ref{fig:jaccard}).
The most stable are geographically compact blocs: the European
(Rotterdam hub, $J=0.602$), Scandinavian (Canada hub, $J=0.760$),
and Pan-American (Panama hub, $J=0.510$) clusters.
The remaining 14 communities dissolved or split, often reflecting
the fragmentation of larger regional blocs into smaller, more
specialised trade clusters---another facet of the rising modularity
trend.

% ====================================================================
\section{Motif analysis}
% ====================================================================

Motif analysis quantifies the statistical over- or
under-representation of each of the 16 directed triad isomorphism
classes \citep{Holland1970} relative to a degree-preserving null
model (100 randomised networks via directed edge swapping).
Z-scores $Z=(f_\text{real}-\bar{f}_\text{rand})/\sigma_\text{rand}$
measure significance.

\begin{figure*}[t]
  \centering
  \includegraphics[width=\linewidth]{Figures/05_motif_zscores.png}
  \caption{Motif Z-scores for all 15 non-null triad classes
    \citep{Holland1970}.
    2015 = blue; 2025 = orange.
    Dashed lines: $Z=\pm2$ significance thresholds.}
  \label{fig:motifs}
\end{figure*}

\begin{figure*}[t]
  \centering
  \includegraphics[width=\linewidth]{Figures/05_motif_paper13_zscores.png}
  \caption{Zoomed comparison for the 13 triad types highlighted
    in \citet{Kaluza2010} figure~4.
    The qualitative pattern closely matches the original study.}
  \label{fig:motifpaper}
\end{figure*}

\begin{figure*}[t]
  \centering
  \includegraphics[width=\linewidth]{Figures/05_motif_frequencies.png}
  \caption{Absolute motif counts for 2015 and 2025.
    ``012'' (no edges) and ``102'' (mutual pair) dominate due
    to network sparsity; ``300'' (complete triad) is small in
    absolute terms but enormously over-represented.}
  \label{fig:motiffreq}
\end{figure*}

\begin{table}[t]
\centering
\caption{Selected triadic census Z-scores.
  $\star$ = types in \citet{Kaluza2010} figure~4.}
\label{tab:motifs}
{\small
\begin{tabular}{@{}llccc@{}}
\toprule
\rowcolor{TableH}
\textbf{Triad} & \textbf{Description} &
$\boldsymbol{Z_{15}}$ & $\boldsymbol{Z_{25}}$ & \textbf{Status} \\
\midrule
012            & No edges        & $-$598  & $-$649  & under \\
102            & Mutual pair     & $+$594  & $+$642  & over  \\
021D$^\star$   & Down fan        & $-$298  & $-$279  & under \\
021U$^\star$   & Up fan          & $-$292  & $-$276  & under \\
021C$^\star$   & Chain           & $-$294  & $-$274  & under \\
111D$^\star$   & Mixed           & $+$131  & $+$116  & over  \\
111U$^\star$   & Mixed (flip)    & $+$128  & $+$119  & over  \\
030T$^\star$   & Triangle (T)    & $-$105  & $-$78   & under \\
030C$^\star$   & Cycle           & $-$114  & $-$92   & under \\
201$^\star$    & 2 mutuals       & $+$373  & $+$404  & over  \\
120D$^\star$   & In-star mut.    & $+$181  & $+$205  & over  \\
120U$^\star$   & Out-star mut.   & $+$147  & $+$149  & over  \\
120C$^\star$   & Chain mutual    & $+$99   & $+$108  & over  \\
210$^\star$    & 5 arcs          & $+$439  & $+$457  & over  \\
300$^\star$    & Complete triad  & $+$1081 & $+$1252 & over  \\
\bottomrule
\end{tabular}}
\end{table}

Figures~\ref{fig:motifs}--\ref{fig:motiffreq} and
table~\ref{tab:motifs} together replicate and extend the motif
analysis of \citet{Kaluza2010} figure~4.
The complete triad (300) has by far the largest Z-score
($+1081$ in 2015, $+1252$ in 2025), confirming that
tightly-connected triangles among hub ports are a defining structural
motif of the GCSN.
Sparse patterns (021D/U/C: fan and chain triads) are strongly
under-represented, consistent with the rich-club structure where
hubs prefer to connect directly to each other rather than through
intermediaries.
The slight increase in $|Z|$ for most classes between 2015 and
2025 is consistent with rising modularity: tighter communities
generate more local cliques (increasing over-representation of
300/210/201) while inter-community fan patterns become even rarer.
The qualitative profile (figure~\ref{fig:motifpaper}) matches the
original study almost exactly, providing strong cross-source
validation.

% ====================================================================
\section{Country-level analysis}
% ====================================================================

\subsection{Network construction and overview}

The three-letter ISO prefix of each GFW port ID
(e.g.\ \texttt{chn-shanghai}$\to$CHN) is parsed to aggregate the
port-level edge list to a country-level network.
Self-loops (intra-country voyages) are removed.

\begin{table}[t]
\centering
\caption{Country-level network summary.}
\label{tab:country}
{\small
\begin{tabular}{@{}lccc@{}}
\toprule
\rowcolor{TableH}
\textbf{Metric} & \textbf{2015} & \textbf{2025} & \textbf{Change} \\
\midrule
Countries (nodes) & 201    & 203    & $+$1.0\%  \\
Directed edges    & 7\,417 & 7\,431 & $+$0.2\%  \\
Mean degree       & 45.68  & 45.43  & ---       \\
Reciprocity       & 0.762  & 0.759  & ---       \\
Clustering $C$    & 0.703  & 0.707  & ---       \\
GCC size          & 100\%  & 100\%  & ---       \\
Communities       & 8      & 7      & $-$1      \\
Modularity $Q$    & 0.566  & 0.586  & $+$3.5\%  \\
\bottomrule
\end{tabular}}
\end{table}

The country-level network (table~\ref{tab:country}) is extremely
dense (mean degree 45.7) and nearly fully bidirectional
(reciprocity 0.76): almost every country pair that trades in one
direction also trades in the reverse.
The GCC covers 100\% of countries---all trading nations are
connected in a single component.
Higher clustering ($C\approx0.70$ vs.\ $0.52$ at port level)
reflects tight regional trade blocs at the national scale.

\subsection{Top countries by shipping strength}

\begin{figure}[t]
  \centering
  \includegraphics[width=\linewidth]{Figures/06_top_countries_strength.png}
  \caption{Top-20 countries by total shipping strength for 2015
    (top) and 2025 (bottom).
    China dominates by a large margin; Netherlands, Belgium, Japan,
    and South Korea follow.}
  \label{fig:topcountries}
\end{figure}

China's dominance is overwhelming in both years
(figure~\ref{fig:topcountries}), with a strength more than twice
that of the second-ranked Netherlands.
The top-5 rankings are stable---Netherlands, Belgium, Japan, South
Korea---reflecting the entrenched positions of major hub-port
nations.
The ranking of the United States falls slightly relative to
Southeast Asian nations, consistent with the US--China trade
tensions redistributing some trans-Pacific volume to intra-Asian
routes.

\subsection{Country strength and degree distributions}

\begin{figure}[t]
  \centering
  \includegraphics[width=\linewidth]{Figures/06_country_distributions.png}
  \caption{Country-level strength and degree distributions
    (log--log) for 2015 and 2025.
    Both are heavy-tailed; the distributions are nearly identical
    between years, confirming structural stability.}
  \label{fig:countrydistr}
\end{figure}

Country-level distributions are heavy-tailed
(figure~\ref{fig:countrydistr}) with a small number of dominant
trading nations.
The near-identical shapes between 2015 and 2025 at the country
level---in contrast to the significant shifts at the port level
(figure~\ref{fig:distr})---suggests that national trade hierarchies
are structurally more stable than individual port rankings.

\subsection{Temporal change at country level}

\begin{figure}[t]
  \centering
  \includegraphics[width=\linewidth]{Figures/06_country_temporal_change.png}
  \caption{Country shipping strength change 2015--2025.
    \textit{Left}: 2015 vs.\ 2025 scatter (log--log);
    diagonal = no change.
    \textit{Right}: Countries with the largest percentage
    gain (green) or loss (red).}
  \label{fig:countrychange}
\end{figure}

Vietnam, India, and Gulf states show the largest percentage gains
(figure~\ref{fig:countrychange}), consistent with industrial
growth and diversification in these economies.
Several legacy East Asian exporters show relative decline,
reflecting production shifts within the region.
The log-log scatter's tight clustering around the diagonal for
large-strength countries confirms that major trading nations
maintained their positions; the most volatile points are small
trading nations with low base volumes.

\subsection{Country flow balance}

\begin{figure*}[t]
  \centering
  \begin{subfigure}[b]{0.28\textwidth}
    \includegraphics[height=0.72\textheight, keepaspectratio]{Figures/06_country_flow_balance_2015.png}
    \caption{2015}
  \end{subfigure}
  \hspace{1.8cm}
  \begin{subfigure}[b]{0.28\textwidth}
    \includegraphics[height=0.72\textheight, keepaspectratio]{Figures/06_country_flow_balance_2025.png}
    \caption{2025}
  \end{subfigure}
  \caption{Country flow balance: net importer vs.\ exporter
    classification in (\textit{a})~2015 and (\textit{b})~2025.
    Most large nations are near-balanced transshipment hubs.
    Asymmetric outliers are typically commodity-specialist terminals
    (bulk mineral exporters or energy import facilities).}
  \label{fig:countrybalance}
\end{figure*}

Most large countries are near-balanced transshipment nodes
(figure~\ref{fig:countrybalance}).
Pronounced asymmetries appear for commodity-specialist economies:
net exporters include bulk mineral and energy producers
(e.g.\ Australia, Saudi Arabia), while net importers include
manufacturing-heavy import consumers.
The flow-balance map is highly stable between years, confirming
that national trade roles---like port rankings---are structurally
locked in at decadal timescales.

\subsection{Country-level communities}

\begin{figure}[t]
  \centering
  \includegraphics[width=\linewidth]{Figures/06_country_communities.png}
  \caption{Louvain trade-bloc communities at country level:
    2015 (8 blocs, $Q=0.566$) and 2025 (7 blocs, $Q=0.586$).
    Communities map to geopolitical regions: East Asia,
    Europe--Mediterranean, Americas, Persian Gulf, Oceania,
    Southeast Asia, and Africa.}
  \label{fig:countrycomm}
\end{figure}

The seven communities in 2025 (reduced from 8, with higher $Q$)
correspond to recognisable geopolitical regions
(figure~\ref{fig:countrycomm}).
The consolidation from 8 to 7 blocs reflects integration of
adjacent regions---likely the absorption of the South Asian bloc
into the broader East Asian trade community following increased
Indian Ocean container traffic.
The higher modularity at the country level ($Q=0.57$--$0.59$)
vs.\ port level ($Q=0.70$--$0.72$) appears paradoxical but reflects
the difference in scale: at the port level, tight local clusters
dominate; at the country level, broader geographic blocs emerge.

\subsection{Country PageRank and centrality}

\begin{figure}[t]
  \centering
  \includegraphics[width=\linewidth]{Figures/06_country_pagerank.png}
  \caption{Top-20 countries by weighted PageRank ($\alpha=0.85$)
    for 2015 and 2025.
    China, Netherlands, Belgium, Japan, and Germany lead.
    Germany and UK rank higher in PR than raw strength,
    reflecting their European gateway roles.}
  \label{fig:countrypr}
\end{figure}

The PageRank ranking (figure~\ref{fig:countrypr}) reveals European
gateway nations---Germany, UK, France---ranking disproportionately
high relative to their raw shipping volumes.
This reflects their role as redistribution hubs: they receive
large container volumes from deep-sea lines and re-export via
feeder services to smaller European and Mediterranean ports.
China's PageRank dominance further confirms that it is the single
most important node in the global trading network by every measure.

\subsection{Trade concentration}

\begin{figure}[t]
  \centering
  \includegraphics[width=\linewidth]{Figures/06_country_lorenz.png}
  \caption{Lorenz curves for country-level shipping strength in
    2015 and 2025.
    Country Gini is substantially lower than port Gini
    (figure~\ref{fig:lorenz}); a mild increase from 2015 to 2025
    signals growing nation-level inequality.}
  \label{fig:countrylorenz}
\end{figure}

Country-level trade concentration (figure~\ref{fig:countrylorenz})
is substantially less extreme than port-level concentration
(figure~\ref{fig:lorenz}).
At the country level, the top 20\% of nations account for
$\sim70\%$ of traffic---high but far from the port-level
concentration where 5\% of routes carry the majority.
The mild Gini increase from 2015 to 2025 at the national level
indicates that the traffic-concentration paradox observed at the
port scale partially aggregates away when viewed nationally.

% ====================================================================
\section{Discussion}
% ====================================================================

\subsection{Replication fidelity}

Our results replicate all qualitative conclusions of
\citet{Kaluza2010}: exponential degree distribution, small-world
clustering, superlinear strength--degree scaling, geographically
coherent modular communities, and over-representation of
clique-like triadic motifs.
Quantitative differences (longer $L$, larger $\mu$ and $\eta$) are
entirely attributable to the larger port universe (7\,097 vs.\ 951
ports) and cross-type aggregation.
The strength--degree scaling exponent $\alpha=1.346$ (2025) is
almost identical to the paper's $\alpha=1.36$, providing strong
cross-source validation.

\subsection{Traffic-concentration paradox}

The central 2015--2025 finding is a \textit{traffic-concentration
paradox}: $+9.7\%$ voyages on $-3.3\%$ routes, with mean degree
$-12.0\%$ and max link weight $+24.4\%$.
This is entirely consistent with the post-2015 deployment of
ultra-large container vessels (ULCVs, $>$18\,000 TEU), which
require deep-water hub ports and generate dramatically higher
per-call voyage counts
\citep{Notteboom2012,Liu2025}.
The structural consequence is a network that is simultaneously more
\textit{efficient per route} (higher $\langle J\rangle$: $9.5\to10.8$)
but less \textit{robust} (higher APL, lower clustering).

\subsection{Growing modularity}

Rising modularity ($Q: 0.699\to0.720$, community count $28\to34$)
indicates tightening regional trade blocs.
This is consistent with US--China trade tensions (post-2018 tariffs)
and pandemic-induced near-shoring, both of which promote
intra-bloc relative to inter-bloc traffic.
The 50\% community stability confirms a stable geographic core
that evolves only gradually, while 14 communities dissolved or
fragmented---likely smaller regional sub-structures responding to
trade-policy shifts.

\subsection{Limitations}

No per-vessel-type breakdown is possible because GFW aggregates
\texttt{CARGO} and \texttt{CONTAINER} types.
Only two snapshots are available; a year-by-year panel would
enable event-based change-point detection (COVID-19 shock,
Suez Canal blockage).
Betweenness centrality and APL estimates are approximations via
random sampling.
No geographic coordinates are available in GFW port IDs,
precluding spatial map visualisation and gravity-model validation.

% ====================================================================
\section{Conclusion}
% ====================================================================

We have presented a comprehensive replication and extension of
\citeauthor{Kaluza2010}'s global cargo ship network analysis using
GFW open data for 2015 and 2025.
Key contributions:
\begin{enumerate}
  \item \textit{Full replication} of Kaluza et~al.\ Table~1 and
        figure~2 on an independent, larger dataset.
  \item \textit{Traffic-concentration paradox}: $+9.7\%$ voyages on
        $-3.3\%$ routes, driven by mega-vessel deployment.
  \item \textit{Centrality extensions}: PageRank identifies East
        Asian gateway ports; k-core ($k_\text{max}=74$) shows
        East Asian ports constitute the densest backbone.
  \item \textit{Chokepoint resilience}: Shanghai is the most
        critical single node; no port removal fragments the GCC.
  \item \textit{50\% community stability} across the decade; rising
        modularity indicates tighter regional blocs.
  \item \textit{Country-level aggregation}: China dominates the
        global trade network; 8 trade blocs consolidate to 7.
  \item \textit{Port specialisation}: large hubs are balanced
        transshipment relays; asymmetric ports are commodity-specialist
        terminals.
\end{enumerate}
All source code and processed edge lists are available at the
project repository.

% ====================================================================
\section*{Acknowledgements}
% ====================================================================

The authors thank the Global Fishing Watch project for open access
to vessel-tracking data.
Analysis used Python~3 with NetworkX \citep{Hagberg2008}, pandas,
NumPy, SciPy, and Matplotlib.

% ====================================================================
\bibliographystyle{plainnat}
\bibliography{references}
% ====================================================================

\end{document}
